%
% 000-session17.tex
% Presentation Session 17: Practice 6 -- Seq2 Server (HTTP Forms)
%
% Compile to .pdf with LaTeX (pdflatex)
% Requires Beamer package (latex-beamer)
%

\documentclass[xcolor=table]{beamer}
\usetheme{Warsaw}
\beamertemplatenavigationsymbolsempty
\setbeamertemplate{headline}{}
\setbeamerfont{title}{size=\large}
\useoutertheme{infolines}
\usepackage[english]{babel}
\usepackage[utf8]{inputenc}
\usepackage{graphics}
\usepackage{amssymb}
\usepackage{multirow}
\usepackage{xcolor}
\usepackage{framed}
\usepackage{hyperref}
\usepackage{tikz}
\usetikzlibrary{arrows.meta, positioning}

\definecolor{shadecolor}{RGB}{180,180,180}

%%%%%%%%%%%%%%%%%%%%%%%%%%%%%%%%%%%%%%%%%%%%%%%%%%%%%%%%%%%%%%%%
% Python Highlighting
%%%%%%%%%%%%%%%%%%%%%%%%%%%%%%%%%%%%%%%%%%%%%%%%%%%%%%%%%%%%%%%%
\DeclareFixedFont{\ttb}{T1}{txtt}{bx}{n}{10}
\DeclareFixedFont{\ttm}{T1}{txtt}{m}{n}{10}

% Custom colors
\usepackage{color}
\definecolor{deepblue}{rgb}{0,0,0.5}
\definecolor{deepred}{rgb}{0.6,0,0}
\definecolor{deepgreen}{rgb}{0,0.5,0}
\definecolor{grey1}{rgb}{0.5,0.5,0.5}

\usepackage{listings}

% Python style for highlighting
\newcommand\pythonstyle{\lstset{
language=Python,
basicstyle=\ttm\small,
morekeywords={self},
keywordstyle=\ttb\color{deepblue},
emphstyle=\ttb\color{deepred},
stringstyle=\color{deepgreen},
commentstyle=\small\color{grey1}\ttm,
frame=single,
showstringspaces=false
}}

% HTML style for highlighting
\newcommand\htmlstyle{\lstset{
language=HTML,
basicstyle=\ttm\small,
keywordstyle=\ttb\color{deepblue},
stringstyle=\color{deepgreen},
commentstyle=\small\color{grey1}\ttm,
frame=single,
showstringspaces=false
}}

% Python environment
\lstnewenvironment{python}[1][]
{
\pythonstyle
\lstset{#1}
}
{}

% HTML environment
\lstnewenvironment{htmlcode}[1][]
{
\htmlstyle
\lstset{#1}
}
{}

% Python for inline
\newcommand\pythoninline[1]{{\pythonstyle\lstinline!#1!}}

% HTML for inline
\newcommand\htmlinline[1]{{\htmlstyle\lstinline!#1!}}

\newcommand{\asignatura}{Session 17: Practice 6 -- Seq2 Server}
\newcommand{\grado}{Programming in Network Environments}
\newcommand{\curso}{2025-2026}

\title[\asignatura]{\asignatura}
\subtitle{\grado}
\author{URJC}
\institute{GSyC}
\date{Course \curso}

\AtBeginSection[]
{
\begin{frame}<beamer>
\begin{center}
{\LARGE \insertsection}
\end{center}
\end{frame}
}

\begin{document}

\frame{
\maketitle
}

\begin{frame}
\frametitle{Contents}
\tableofcontents
\end{frame}

%%%%%%%%%%%%%%%%%%%%%%%%%%%%%%%%%%%%%%%%%%%%%%%%%%%%%%%%%%%%%%%%
%%%%%%%%%%%%%%%%%%%%%%%%%%%%%%%%%%%%%%%%%%%%%%%%%%%%%%%%%%%%%%%%
% SECTION 1: Session 16 Review
%%%%%%%%%%%%%%%%%%%%%%%%%%%%%%%%%%%%%%%%%%%%%%%%%%%%%%%%%%%%%%%%
%%%%%%%%%%%%%%%%%%%%%%%%%%%%%%%%%%%%%%%%%%%%%%%%%%%%%%%%%%%%%%%%

\section{Review: Session 16 -- HTML Forms}

\begin{frame}[fragile]
\frametitle{What We Learned in Session 16 (1/2)}

\textbf{HTML Forms}: allow clients to \textbf{send data} to web servers.

\begin{itemize}
\item \textbf{Form structure}: \htmlinline{<form action="..." method="get">}
  \begin{itemize}
  \item \texttt{action}: the server \textbf{endpoint} to send data to
  \item \texttt{method}: HTTP method (\textbf{GET} puts data in the URL)
  \end{itemize}
\vspace{0.15cm}
\item \textbf{Input types}: text, checkbox, radio, dropdown (select)
\vspace{0.15cm}
\item \textbf{Query string}: data is appended to the URL after \texttt{?}
  \begin{itemize}
  \item \texttt{/echo?msg=Hello\&chk=on} -- key=value pairs separated by \texttt{\&}
  \end{itemize}
\end{itemize}

\end{frame}

\begin{frame}[fragile]
\frametitle{What We Learned in Session 16 (2/2)}

\textbf{Parsing form data} on the server side:

\vspace{0.1cm}

\begin{python}
from urllib.parse import parse_qs, urlparse

# Inside do_GET():
url_path = urlparse(self.path)
path = url_path.path       # "/echo"
arguments = parse_qs(url_path.query)
# {"msg": ["Hello"]}
\end{python}

\vspace{0.2cm}

\textbf{Dynamic HTML} with \pythoninline{jinja2}:

\begin{itemize}
\item \pythoninline{j.Template(contents).render(context=\{...\})}
\item In HTML: \texttt{\{\{context.todisplay\}\}} for dynamic values
\end{itemize}

\end{frame}

%%%%%%%%%%%%%%%%%%%%%%%%%%%%%%%%%%%%%%%%%%%%%%%%%%%%%%%%%%%%%%%%
%%%%%%%%%%%%%%%%%%%%%%%%%%%%%%%%%%%%%%%%%%%%%%%%%%%%%%%%%%%%%%%%
% SECTION 2: Session 17 -- Practice 6
%%%%%%%%%%%%%%%%%%%%%%%%%%%%%%%%%%%%%%%%%%%%%%%%%%%%%%%%%%%%%%%%
%%%%%%%%%%%%%%%%%%%%%%%%%%%%%%%%%%%%%%%%%%%%%%%%%%%%%%%%%%%%%%%%

\section{Session 17: Practice 6 -- Seq2 Server}

\begin{frame}
\frametitle{Session 17 Goals}

\textbf{In this session we will}:

\begin{itemize}
\item Develop a second version of the \textbf{Seq server} from Practice 3
\item Use the \textbf{HTTP Python module} instead of raw sockets
\item Access services through \textbf{HTML forms}
\item Implement 4 services: \textbf{PING}, \textbf{GET}, \textbf{GENE}, \textbf{OPERATION}
\end{itemize}

\vspace{0.4cm}

\textbf{Time}: 2 hours

\vspace{0.2cm}

\textbf{Folder}: \texttt{P06/} with subfolder \texttt{html/}

\vspace{0.2cm}

\textbf{Important}: Finish all S16 exercises before starting!

\end{frame}

\begin{frame}
\frametitle{Services Overview}

\begingroup\small
\begin{center}
\begin{tabular}{|l|l|l|p{3cm}|}
\hline
\textbf{Service} & \textbf{Action} & \textbf{Argument} & \textbf{Description} \\
\hline
PING & \texttt{ping} & none & Test if server is alive \\
\hline
GET & \texttt{get} & \texttt{n}: 0--4 & Get sequence number n \\
\hline
GENE & \texttt{gene} & \texttt{name} & Get complete gene sequence \\
\hline
OPERATION & \texttt{operation} & \texttt{seq}, \texttt{op} & Perform operation on sequence \\
\hline
\end{tabular}
\end{center}
\endgroup

\vspace{0.2cm}

\textbf{Valid gene names}: U5, ADA, FRAT1, FXN, RNU6\_269P

\vspace{0.2cm}

\textbf{Valid operations}: \texttt{info}, \texttt{comp} (complement), \texttt{rev} (reverse)

\end{frame}

%%%%%%%%%%%%%%%%%%%%%%%%%%%%%%%%%%%%%%%%%%%%%%%%%%%%%%%%%%%%%%%%
%%%%%%%%%%%%%%%%%%%%%%%%%%%%%%%%%%%%%%%%%%%%%%%%%%%%%%%%%%%%%%%%
% SECTION 3: Exercise 1 -- PING
%%%%%%%%%%%%%%%%%%%%%%%%%%%%%%%%%%%%%%%%%%%%%%%%%%%%%%%%%%%%%%%%
%%%%%%%%%%%%%%%%%%%%%%%%%%%%%%%%%%%%%%%%%%%%%%%%%%%%%%%%%%%%%%%%

\section{Exercise 1: PING}

\begin{frame}[fragile]
\frametitle{Exercise 1: PING Service}

\textbf{Implement the PING service} in the server.

\vspace{0.2cm}

\begin{itemize}
\item Request \texttt{/} $\rightarrow$ return the \textbf{main page} (\texttt{index.html}) with forms
\item Form submits to \texttt{/ping} $\rightarrow$ server responds with a page saying the \textbf{server is alive}
\item Include a \textbf{link} back to the main page
\item Any other resource $\rightarrow$ return an \textbf{error page}
\end{itemize}

\vspace{0.3cm}

\textbf{Files}:
\begin{itemize}
\item \texttt{P06/Seq2-server.py} -- The server
\item \texttt{P06/html/index.html} -- Main page with forms
\item \texttt{P06/html/error.html} -- Error page
\end{itemize}

\end{frame}

\begin{frame}[fragile]
\frametitle{Server Structure for Seq2}

\begin{python}
import http.server
import socketserver
from urllib.parse import parse_qs, urlparse

PORT = 8080
socketserver.TCPServer.allow_reuse_address = True

class Seq2Handler(
        http.server.BaseHTTPRequestHandler):

    def do_GET(self):
        url_path = urlparse(self.path)
        path = url_path.path
        arguments = parse_qs(url_path.query)

        # HERE, PROCESS url_path AND arguments
\end{python}

\end{frame}

%%%%%%%%%%%%%%%%%%%%%%%%%%%%%%%%%%%%%%%%%%%%%%%%%%%%%%%%%%%%%%%%
%%%%%%%%%%%%%%%%%%%%%%%%%%%%%%%%%%%%%%%%%%%%%%%%%%%%%%%%%%%%%%%%
% SECTION 4: Exercise 2 -- GET
%%%%%%%%%%%%%%%%%%%%%%%%%%%%%%%%%%%%%%%%%%%%%%%%%%%%%%%%%%%%%%%%
%%%%%%%%%%%%%%%%%%%%%%%%%%%%%%%%%%%%%%%%%%%%%%%%%%%%%%%%%%%%%%%%

\section{Exercise 2: GET}

\begin{frame}
\frametitle{Exercise 2: GET Service}

\textbf{Extend the server} to implement the GET service.

\vspace{0.2cm}

\begin{itemize}
\item The form sends a \textbf{sequence number} (0--4) to the server
\item The server responds with the \textbf{requested sequence}
\item Include a \textbf{link} back to the main page
\end{itemize}

\vspace{0.3cm}

\textbf{New file}: \texttt{P06/html/get.html} -- Form for requesting a sequence

\vspace{0.3cm}

\textbf{Example request}:

\vspace{0.1cm}

\begin{center}
\texttt{GET /get?n=2 HTTP/1.1}
\end{center}

\vspace{0.2cm}

The server should have 5 predefined sequences (numbered 0 to 4) and return the one requested.

\end{frame}

%%%%%%%%%%%%%%%%%%%%%%%%%%%%%%%%%%%%%%%%%%%%%%%%%%%%%%%%%%%%%%%%
%%%%%%%%%%%%%%%%%%%%%%%%%%%%%%%%%%%%%%%%%%%%%%%%%%%%%%%%%%%%%%%%
% SECTION 5: Exercise 3 -- GENE
%%%%%%%%%%%%%%%%%%%%%%%%%%%%%%%%%%%%%%%%%%%%%%%%%%%%%%%%%%%%%%%%
%%%%%%%%%%%%%%%%%%%%%%%%%%%%%%%%%%%%%%%%%%%%%%%%%%%%%%%%%%%%%%%%

\section{Exercise 3: GENE}

\begin{frame}
\frametitle{Exercise 3: GENE Service}

\textbf{Extend the server} to implement the GENE service.

\vspace{0.2cm}

\begin{itemize}
\item The form sends a \textbf{gene name} to the server
\item The server responds with the \textbf{complete sequence} of the given gene
\item Include a \textbf{link} back to the main page
\end{itemize}

\vspace{0.3cm}

\textbf{New file}: \texttt{P06/html/gene.html} -- Form for requesting a gene

\vspace{0.3cm}

\textbf{Valid gene names}: U5, ADA, FRAT1, FXN, RNU6\_269P

\vspace{0.3cm}

\textbf{Example request}:

\vspace{0.1cm}

\begin{center}
\texttt{GET /gene?name=U5 HTTP/1.1}
\end{center}

\end{frame}

%%%%%%%%%%%%%%%%%%%%%%%%%%%%%%%%%%%%%%%%%%%%%%%%%%%%%%%%%%%%%%%%
%%%%%%%%%%%%%%%%%%%%%%%%%%%%%%%%%%%%%%%%%%%%%%%%%%%%%%%%%%%%%%%%
% SECTION 6: Exercise 4 -- OPERATION
%%%%%%%%%%%%%%%%%%%%%%%%%%%%%%%%%%%%%%%%%%%%%%%%%%%%%%%%%%%%%%%%
%%%%%%%%%%%%%%%%%%%%%%%%%%%%%%%%%%%%%%%%%%%%%%%%%%%%%%%%%%%%%%%%

\section{Exercise 4: OPERATION}

\begin{frame}
\frametitle{Exercise 4: OPERATION Service}

\textbf{Extend the server} to implement the OPERATION service.

\vspace{0.2cm}

Three operations on a given sequence:
\begin{enumerate}
\item \textbf{info}: Get information about the sequence (length, base counts and percentages)
\item \textbf{comp}: Calculate the \textbf{complement}
\item \textbf{rev}: Calculate the \textbf{reverse}
\end{enumerate}

\vspace{0.2cm}

\textbf{New file}: \texttt{P06/html/operation.html} -- Form for operations

\vspace{0.2cm}

\textbf{Example request}:

\vspace{0.1cm}

\begin{center}
\texttt{GET /operation?seq=ATAGAC\&op=info HTTP/1.1}
\end{center}

\end{frame}

\begin{frame}[fragile]
\frametitle{INFO Response Format}

The \textbf{INFO} operation should return a response like this:

\vspace{0.2cm}

\begin{python}[language={}]
Sequence: ATAGACCAAACATGAGAGGCT
Total length: 21
A: 9 (42.9%)
C: 4 (19.0%)
G: 5 (23.8%)
T: 3 (14.3%)
\end{python}

\vspace{0.2cm}

\begin{itemize}
\item Count each \textbf{base} (A, C, G, T) in the sequence
\item Calculate the \textbf{percentage} of each base
\item Display the \textbf{total length} of the sequence
\end{itemize}

\end{frame}

%%%%%%%%%%%%%%%%%%%%%%%%%%%%%%%%%%%%%%%%%%%%%%%%%%%%%%%%%%%%%%%%
%%%%%%%%%%%%%%%%%%%%%%%%%%%%%%%%%%%%%%%%%%%%%%%%%%%%%%%%%%%%%%%%
% SECTION 7: Deliverables
%%%%%%%%%%%%%%%%%%%%%%%%%%%%%%%%%%%%%%%%%%%%%%%%%%%%%%%%%%%%%%%%
%%%%%%%%%%%%%%%%%%%%%%%%%%%%%%%%%%%%%%%%%%%%%%%%%%%%%%%%%%%%%%%%

\section{Deliverables}

\begin{frame}
\frametitle{Deliverables}

\textbf{P06 folder should contain}:

\vspace{0.2cm}

\begingroup\small
\begin{itemize}
\item \texttt{Seq2-server.py} -- The Seq2 web server
\item \texttt{html/index.html} -- Main page with forms for all services
\item \texttt{html/ping.html} -- Ping response page
\item \texttt{html/get.html} -- Get form page
\item \texttt{html/gene.html} -- Gene form page
\item \texttt{html/operation.html} -- Operation form page
\item \texttt{html/error.html} -- Error page
\end{itemize}
\endgroup

\vspace{0.3cm}

\textbf{Remember}: Push everything to your remote Github repository!

\end{frame}

%%%%%%%%%%%%%%%%%%%%%%%%%%%%%%%%%%%%%%%%%%%%%%%%%%%%%%%%%%%%%%%%
%%%%%%%%%%%%%%%%%%%%%%%%%%%%%%%%%%%%%%%%%%%%%%%%%%%%%%%%%%%%%%%%
% SECTION 8: End
%%%%%%%%%%%%%%%%%%%%%%%%%%%%%%%%%%%%%%%%%%%%%%%%%%%%%%%%%%%%%%%%
%%%%%%%%%%%%%%%%%%%%%%%%%%%%%%%%%%%%%%%%%%%%%%%%%%%%%%%%%%%%%%%%

\begin{frame}
\begin{center}
\Huge{Questions?}

\vspace{1cm}

\Large{Time to code!}
\end{center}
\end{frame}

\end{document}
